\begin{frame}{Symmetrische vs. Asymmetrische Kryptosysteme}
	\begin{block}{Symmetrische Kryptosysteme (z.B. Caesar, Vigenère, OTP)}
		Verwenden \textbf{demselben Schlüssel} zur Ver- und Entschlüsselung. Problem: Wie tauscht man den Schlüssel sicher aus?
	\end{block}

	\pause
	\begin{myidea}[Asymmetrische Kryptosysteme (Public-Key Kryptologie)]
		Jede Person generiert ein Schlüsselpaar:
		\begin{itemize}[<+->]
			\item \textbf{Öffentlicher Schlüssel} (\emoji{key}): Wird allen bekannt gemacht. Jeder kann damit Nachrichten für den Empfänger verschlüsseln.
			\item \textbf{Privater Schlüssel} (\emoji{lock}): Bleibt absolut geheim. Nur der Empfänger kann damit die Nachrichten entschlüsseln.
		\end{itemize}
	\end{myidea}
\end{frame}

\begin{frame}{Prinzip der Public-Key-Verschlüsselung}
	\begin{figure}[H]
		\centering
		\resizebox{!}{0.75\textheight}{\begin{tikzpicture}[
	node distance = 0.9cm and 2.0cm,
	data/.style={rectangle, rounded corners, text width=2.8cm, minimum height=0.95cm, text centered, draw=ForestGreen, fill=ForestGreen!20, font=\small},
	keybox/.style={rectangle, rounded corners, text width=2.8cm, minimum height=0.95cm, text centered, draw=RoyalBlue, fill=RoyalBlue!20, font=\small},
	lbl/.style={font=\footnotesize\color{black!80}}
]
	% Nodes
	\node (plaintext) [data] {\emoji{woman-technologist} Alice};
	\node (ciphertext) [data, right=of plaintext] {Verschlüsselter\\Text};
	\node (publickey) [keybox, above=of ciphertext] {\emoji{key} Öffentlicher Schlüssel\\von Bob};

	\node (decryptedtext) [data, below=of ciphertext] {Verschlüsselter\\Text};
	\node (privatekey) [keybox, below=of decryptedtext] {\emoji{lock} Privater Schlüssel\\von Bob};
	\node (bob) [data, left=of decryptedtext] {\emoji{man-technologist} Bob};

	% Arrows
	\draw [->, thick] (plaintext) -- (ciphertext);
	\draw [->, thick] (ciphertext) -- (decryptedtext);
	\draw [->, thick] (publickey) -- (ciphertext);
	\draw [->, thick] (privatekey) -- (decryptedtext);
	\draw [->, thick] (decryptedtext) -- (bob);

	% Labels
	\node at ($(plaintext.east)!0.5!(ciphertext.west)$) [above, lbl] {Klartext};
	\node at ($(publickey.south)!0.5!(ciphertext.north)$) [left, lbl] {Verschlüsselung};
	\node at ($(ciphertext.south)!0.5!(decryptedtext.north)$) [left, lbl] {Übertragung};
	\node at ($(privatekey.north)!0.5!(decryptedtext.south)$) [left, lbl] {Entschlüsselung};
	\node at ($(decryptedtext.west)!0.5!(bob.east)$) [above, lbl] {Klartext};
\end{tikzpicture}
}
		\caption{Prinzip der Public-Key-Verschlüsselung}
	\end{figure}
\end{frame}

\begin{frame}{Das RSA-Verfahren}
	1977 von Ron Rivest, Adi Shamir und Leonard Adleman erfunden.

	\begin{mydefinition}[Mathematisches Fundament von RSA]
		Es basiert auf der Schwierigkeit der \textbf{Primfaktorzerlegung}:
		\begin{itemize}
			\item Zwei grosse Primzahlen $p$ und $q$ miteinander zu multiplizieren ($n = p \cdot q$) ist extrem einfach.
			\item Aus dem Produkt $n$ wieder die Primfaktoren $p$ und $q$ zu finden, ist für ausreichend grosse Zahlen praktisch unmöglich!
		\end{itemize}
	\end{mydefinition}
\end{frame}

\begin{frame}{RSA Schlüsselerzeugung: Theorie \& Beispiel}
	\begin{columns}[T]
		\begin{column}{0.50\textwidth}
			\begin{block}{Theorie / Algorithmus}
				{\scriptsize
					\begin{enumerate}
						\item Wähle zwei Primzahlen $p$ und $q$.
						\item Modul: $n = p \cdot q$
						\item Phi-Wert: $\varphi(n) = (p-1)(q-1)$
						\item Wähle $e$ mit $\operatorname{ggT}(e, \varphi(n)) = 1$.
						\item Berechne $d$ mit:
						      \[ (e \cdot d) \equiv 1 \pmod{\varphi(n)} \]
						\item \textbf{Public}: $(e, n)$, \textbf{Private}: $(d, n)$.
					\end{enumerate}
				}
			\end{block}
		\end{column}
		\begin{column}{0.48\textwidth}
			\begin{exampleblock}{Zahlenbeispiel}
				{\scriptsize
					\begin{itemize}
						\item Wähle $p = 3, \; q = 11$.
						\item Modul: $n = 3 \cdot 11 = 33$
						\item $\varphi(33) = (3-1)(11-1) = 20$
						\item Wähle $e = 3$ (da $\operatorname{ggT}(3, 20) = 1$).
						\item Bestimme $d$:
						      \[ 3 \cdot d \equiv 1 \pmod{20} \implies d = 7 \]
						\item \textbf{Public}: $(3, 33)$, \textbf{Private}: $(7, 33)$.
					\end{itemize}
				}
			\end{exampleblock}
		\end{column}
	\end{columns}
\end{frame}

\begin{frame}{RSA Ver- und Entschlüsselung: Theorie \& Beispiel}
	\begin{columns}[T]
		\begin{column}{0.50\textwidth}
			\begin{block}{Theorie}
				{\scriptsize
					Public Key $(e, n)$, Private Key $(d, n)$.\\[0.3em]
					\textbf{1. Verschlüsselung (Absender)}:
					\begin{itemize}
						\item Klartext als Zahl $m < n$.
						\item Chiffrat berechnen:
						      \[ c = m^e \bmod n \]
					\end{itemize}
					\textbf{2. Entschlüsselung (Empfänger)}:
					\begin{itemize}
						\item Mit privatem Schlüssel $d$:
						      \[ m = c^d \bmod n \]
					\end{itemize}
					\textbf{Garantie}: $m \equiv (m^e)^d \pmod n$ (Euler).
				}
			\end{block}
		\end{column}
		\begin{column}{0.48\textwidth}
			\begin{exampleblock}{Zahlenbeispiel}
				{\scriptsize
					\textbf{Schlüssel}: Public $(3, 33)$, Private $(7, 33)$.\\[0.3em]
					\textbf{1. Verschlüsselung} von Klartext $m = 3$:
					\[ c = 3^3 \bmod 33 = 27 \bmod 33 = \mathbf{27} \]
					\textbf{2. Entschlüsselung} von Chiffrat $c = 27$:
					\[ m = 27^7 \bmod 33 = \mathbf{3} \quad \textcolor{ForestGreen}{\mathbf{\checkmark}} \]
				}
			\end{exampleblock}
		\end{column}
	\end{columns}
\end{frame}

\begin{frame}{Auftrag: RSA-Verfahren}{Skript}
	\aufgabebox{Aufgabe 4.1 (RSA Schlüsselerzeugung \& Verschlüsselung)}
	\challengebox{4.3: RSA mit Python für eigene Nachrichten}
\end{frame}
